\documentclass[
  %a4paper,
  12pt]
{scrbook}%reprt} 
\RequirePackage[german, french]{babel} %French is the default
\RequirePackage[T1]{fontenc}
\RequirePackage{scrlayer-scrpage}
\RequirePackage{comment}
\RequirePackage{exercise}
%Workaround
\renewcommand\ExerciseName{Exercice}
\renewcommand\AnswerName{Solution de l'exercice}
\RequirePackage[short, nodayofweek]{datetime}
\RequirePackage{xr} % {xr-hyper} % {zref-xr} % \externaldocument{chapterI}
\externaldocument[O-]{lecondall-original}% [O-]{lecondall-original.tex}
%\usepackage{sidenotes}%{marginnote} 
%\usepackage[useregional]{datetime2}
%\usepackage{paracol}%textcol
\usepackage{graphbox}
%\usepackage[usenames, dvipsnames]{color} % textcol
%\usepackage{yfonts} % textfrak
\usepackage{natbib} %[numbers]
%\usepackage{biblatex}
\usepackage{anyfontsize}
\usepackage{t1enc}
\usepackage{yfonts}
\usepackage{latexsym}
\newcommand\HUGE{\fontsize{50}{60} \selectfont}
%\newcommand\HUGE{\fontsize{50}{60} \selectfont}
\usepackage{hyperref}

\newcommand\germanbegin[1]{\begin{otherlanguage*}{german}}
\newcommand\germanend[1]{\end{otherlanguage*}{german}}
\newcommand\envirgerman[1]{\begin{otherlanguage*}{german}#1\end{otherlanguage*}}
\input{ERWANN-macro-fr-notes.tex}
% Envir

\newcommand\LDAdemarking[1]{#1
%bug prone, so leave empty, for now
%\colorbox{green}{#1}
%%\textcolor{green}{#1} 
}


\newcommand\LDAde[1]{% Warning: chapters don't take envirgerman
	\envirgerman{#1}
}

% Marking
\newcommand\ldaBLANK{\underline{\hspace{1cm}}}


% Structure

\newcommand\ldaTTL{La Le\c con d'Allemand Syst\'ematique}

\newcommand\ldaTXT{Texte de base}
\newcommand\ldaLEX{Rubrique lexicale}
\newcommand\ldaLEXa{$1\geq2$}
\newcommand\ldaLEXb{$10/10$}
\newcommand\ldaGRAM{Rubrique grammaticale}
\newcommand\ldaEX[1]{Exercice #1}
\newcommand\ldaEXs{Exercices}
\newcommand\ldaLEXc{Approfondissement lexical}
\newcommand\ldaTHEME{Th\`eme de contr\^ole}
\newcommand\ldaTHEMEsol{Solution au th\`eme de contr\^ole}
\newcommand\ldaDISK{Diskussions - und Aufsatzthemen}

\newcommand\LDAint{intitul\'e}%what for?

\newcommand\ldaAUTHOR{Jean-Pierre Demarche}

% Notes

%\newcommand*\LDAnote[1]{\sidenote{#1}}
\newcommand*\ldaVERIF{\noteverifier}
\newcommand*\ldaAF{\notefaire}
\newcommand*\ldaALT{\notesinon}

% Remplacements

\newcommand\ldaVERBE[1]{#1}% sur le mod\`ele de
\newcommand\ldaCONJ[1]{#1}%conjonction
\newcommand\ldaFORME[1]{#1}%forme
\newcommand\ldaUSAGE[1]{#1}%usage
\newcommand\ldaPREPOS[1]{#1}%preposition

% Autres

\newcommand\ldaVERBEpart[2]{#1|#2}
%\newcommand\LDAindicsol[2]{\begin{tabular}{l|l}#1&#2\end{tabular}}
\newcommand\ldaSOL[1]{\textbf{\noindent Solution:}}
\newcommand\ldaINDIC{\textbf{\noindent Indication:}}




\pagestyle{scrheadings}
\rohead[Copyright (c) {\censor{Erwann Rogard}}, 2016]{}
\lohead[Recipient:]{}


%\newcommand{\mydate}{\DTMdisplaydate{1994}{6}{3}{-1}}
 
\begin{document}

\begin{titlepage}
%{\HUGE La le\c con d'Allemand syst\'ematique\\ Solutions}
%}
\begin{center}
\includegraphics[
	align=c,
	width=0.5\textwidth]{berlincoatofarms.png}
%  \qquad
\begin{minipage}{0.45\textwidth}
%\textswab{Allemand}
\textbf{\HUGE La le\c con\\ d'\textgoth{Allemand} syst\'ematique.}
\end{minipage}
\vspace*{\fill}

\textbf{\HUGE Les solutions aux exercices}

\vspace*{\fill}

\textbf{{\censor{Erwann Rogard}}; version du \today}

\end{center}
\end{titlepage}

\chapter*{Introduction}

Ce livre contient mes r\'eponses aux exercices (entendus au sens large, explication \`a suivre) de la premi\`ere partie du livre pour les classes pr\'eparatoires et de lyc\'ee intitul\'e \emph{\ldaTTL}\citep{demarche1989}, pour la suite, l'\og original\fg. 

J'ai fait ce travail pour que mon apprentissage de l'Allemand, qui remonte \`a mes ann\'ees de coll\`ege, lyc\'ee (le bac en 1991) et fac, quasiment dormant depuis lors, mais remis \`a jour, un peu par hasard, entre l'automne 2013 et le printemps 2014, ne reste pas en friche. 

Pourquoi \emph{\ldaTTL}? Parce qu'il avait surv\'ecu aux \'elaguements successifs de mes affaires scolaires et universitaires. Avec ses textes ardus et `vintage' \footnote{Recensement des ann\'ees de publication des `\ldaTXT' des 20 le\c cons syst\'ematiques: 1944, 1987, 1983, 1978, 1985, 1981,  1925, 1985, 1944, 1985, 1986, 1977, 1987, 1981, 1981, 1985, 1987, 1986, 1986.}, en arrivant au bout de ses 20 le\c cons syst\'ematiques, j'aurais ainsi le sentiment de la t\^ache accomplie. 

Dans l'original, chaque chapitre est structur\'e comme suit. Figurent en gras les sections pour lesquels ce livre apporte des propositions de solutions: 
\begin{itemize}
\item \ldaTXT
\item \ldaLEX, comprenant deux sous-sections, \ldaLEXa~et \ldaLEXb. 
\item \ldaGRAM
\item \textbf{\ldaEXs}
\item \ldaLEXc
\item \textbf{\nameref{O-spiel:theme}}
\item \nameref{O-spiel:disk} (mini-essais \`a r\'ediger par l'\'el\`eve)
\end{itemize}
%Ici, les deux derni\`eres sections sont int\'egr\'ees \`a la liste des exercices num\'erot\'es, avec en titre leur intitul\'e d'origine pour les distinguer du reste. %Par exemple, au \autoref{O-ch:spiel}, ces deux sections sont \og\mbox{\ExerciseName~\autoref{spiel:theme} \qquad \nameref{O-spiel:theme}}\fg~et \og\mbox{\ExerciseName~\autoref{spiel:disk} \qquad \nameref{O-spiel:disk}}\fg.

Les \'enonc\'es sont repris tels quels, \'eventuellement amend\'es d'un renvoi \`a la section de l'original qui lui est pertinente, soit dans le titre de l'exercice, soit pr\'ec\'ed\'e de `\ldaINDIC', \`a la suite de l'\'enonc\'e. 

Ce livre a \'et\'e r\'ealis\'e avec la le logiciel libre \LaTeX, pr\'ecis\'ement avec la classe \texttt{scrbook}.

%\chapter{\LDAdemarking{\nameref{ch:mann}}}

\section*{\LDAex{1}}

%R�unissez les deux phrases en une seule, la deuxi�me devenant une � principale � en t�te de phrase, la deuxi�me une � infinitive compl�ment � (Faites attention: des suppressions de mots et des transferts d'une proposition � l'autre peuvent �tre n�cessaires!);  

\begin{enumerate}
  
\item%1 Ein Rettungsboot wird auf das Wasser hinuntergelassen. Der Kapit�n befiehlt es.

\LDAde{Der Kapit\"an befiehlt es, ein Rettungsboot auf das Wasser hinunterzulassen.}

\item%2 J�rgen darf heute Abend nicht ins Kino gehen. Seine Eltern haben es ihn verboten.

\LDAde{Die Eltern von J\"urgen haben ihm verboten, heute Abend ins Kino zu gehen.}

\item%3 Der betrunkene Matrose kann nicht herausgefischt werden. Die Mannschaft versucht er vergeblich. 

\LDAde{Die Mannschaft versucht vergeblich den betrunkene Matrose herauszufischen.}

\item%4 Ich tra�me nicht davon. Ich wage es nicht damal.

\LDAde{Ich h\"atte einmal nicht gewagt, davon zu tr\"aumen.}

\item%5 Er repariert sein altes Auto. Ich werde ihm dabei helfen.

\LDAde{Ich werde ihm helfen, sein altes Auto zu reparieren.}

\item%6 Ich werde die Polizei nicht verst�ndigen. Ein Herr hat mir davon abgeraten.

\LDAde{Ein Herr hat mir abgeraten, dir Polizei zu verst\"andigen.}

\item%7 Der einzige �berlebende kann von der Erz�hlung des Schiffbruchs gut leben. Er hat es gehofft.

\LDAde{Der einzige \"Uberlebende hat es gehofft, von der Erz\"ahlung des Schiffbruchs, gut leben zu k\"onnen.}

\item%8 Mein Sohn darf am Samstag abend l�nger aufbleiben. Et hat mich darum gebeten.

\LDAde{Mein Sohn hat mich gebeten, am Samstagabend l\"anger aufbleiben zu d\"urfen.}

\item%9 Er f�hrt auf mich zu. Ich sehe es von weitem.

\LDAde{Ich sehe est von weiten, er auf mich zuf\"ahrt.}

%[zufahren: http://ow.ly/Zo4kg; Jamais de 'zu' apr�s un verbe de perception]

\item%10 Er kann nicht um Hilfe rufen. Er hat nicht mehr die Kraft dazu.

\LDAde{Er hat nicht mehr die Kraft, um Hilfe rufen zu k\"onnen.}

%[Jamais de zu apr�s (mais pas avant!) un auxiliaire de mode]

\item%11 Das Schiff lief auf ein Korallenriff auf. Alle Matrosen h�rten es.

\LDAde{Alle Matrosen h\"orten das Schiff auf ein Korallenriff auflaufen.}

\item%12 Ich lasse mich nicht betr\"ugen. Ich dulde es einfach nicht.

\LDAde{Ich dulde es einfach nicht, mich betr\"ugen zu lassen.}

\item%13 Sie darf etwas sp�ter ankommen. Sie hat den Wunsch ausgedr�ckt.

\LDAde{Sie hat den Wunsch ausgedr\"uckt, etwas sp\"ater ankommen zu d\"urfen.}

%[Um es anders auszudr�cken / Pour le dire autrement]

\item%14 Er will nicht zusehen. Er hat Angst.

\LDAde{Er hat Angst zusehen.}

\item%15 Sie k�nnen die Rettung ihres Freundes bald feiern. Sie freuen sich.

\LDAde{Sie freuen sich, die Rettung ihres Freundes bald feiern zu k\"onnen.}

\end{enumerate}

\section*{\LDAex{2}}

%Traduisez les phrases que vous avez faites dans l'exercice pr�c�dent. 

\begin{enumerate}

\item%1
  Le capitaine a ordonn\'e que le canot de sauvetage soit mis \`a la mer.

\item%2
  Les parents de J\"urgen lui ont interdit d'aller ce soir au cin\'ema.

\item%3
  L'\'equipe essaye sans succ\`es de tirer de l'eau le marin so\^ul.

\item%4
  Je n'aurais m\^eme pas os\'e en r\^ever.

\item%5
  Je vais l'aider \`a r\'eparer sa voiture.

\item%6
  Un monsieur m'a d\'econseill\'e d'informer la police.

\item%7 
  Le seul survivant du naufrage esp\'erait que le r\'ecit de celui-ci lui ferait surmonter l'\'epreuve.

\item%8
  Mon fils ma demand\'e la permission de se coucher plus tard samedi en soir\'ee.

\item%9
  Je vois de loin qu'il me fonce dessus.

\item%10
Il n'a plus la force d'appeler \`a l'aide.  

\item%11
  Tous les matelots entendirent le navire s'\'echouer contre la barri\`ere de corail.

\item%12
  Je ne tol\`ere pas facilement que l'on me trompe.

\item%13
  Elle a fait exprim\'e le souhait d'\^etre autoris\'ee \`a venir plus tard.

\item%14
  Il a peur de regarder.

\item%15
  Ils se r\'ejouissent de pouvoir prochainement f\^eter le sauvetage de leur ami.
  
\end{enumerate}

\section*{\LDAex{3}}

%Mettez les deux premiers paragraphes du texte au pr�sent de l'indicatif (pr�sent narratif!)

\LDAde{Der Wind weht nicht so stark. Bei einem schlingern des Schiffes verliert der Matrose das Gleichgewicht und st\"urzt von Deck. Der Mann am Ruder sieht den Sturz und gibt sofort Alarm. Der Kapit\"an befiehl, ein Boot auf das m\"assig bewegte Wasser hinunterzulassen, den langsam forttreibenden Matrosen zu retten. }

\LDAde{Die Mannschaft liegt sich kr\"aftig in die Riemen, und schon nach wenigen Schl\"agen erreichen sie den um Hilfe Rufenden. Sie werfen ihm einen Rettungsring zu, and der sich klammert. Im n\"aherschaukelnden Boot richtet sich im Bug einer auf, um dem \ERquoteskip, doch verliert der Retter selbst den Halt und in die Fluten f\"allt, w\"ahrend eine ungeahnte Woge das Boot seitlich unter l\"auft und umwirft. Der Kapit\"an gibt Anweisung, \ERquoteskip. Doch kaum hat man damit begonnen, ersch\"uttert ein Stoss das Schiff, das sich schon zur Seite liegt, \ERquoteskip, das sich knapp under der Oberfl\"ache verbirgt. Der Kapit\"an versackt wie \"ublich mit dem t\"odlich verwundeten Schiff.}

\section*{\LDAex{4}}

%Int�grez les verbes qui figurent entre parenth�ses à la suite des phrases suivantes (Attention: la mention «auxiliaire» signifie que vous devez choisir celui qui est ad�quat!); Traduisez les phrases que vous avez faites dans l'exercice pr�c�dent.

\begin{enumerate}

\item% 1 Ich bin der Meinung, er (...) diesem Zwischenfall eine zu grosse Bedeutung (...). (auxiliaire; beimessen)

\LDAde{hat; beigemessen}

\item%2 Frau Steiner ist w\"utend, weil die Kinder die Tafel Schokolade (...) (...). (aufessen; auxiliaire)

\LDAde{aufgegessen; haben}

\item%3 Er m\"usste schon l�ngst hier sein; hoffentlich (...) ihm unterwegs nichts (...). (auxiliaire; geschehen)

\LDAde{w\"urde; geschehen}

\item%4 Wenn sie Zeit hat, (...) sie jeden Tag die Lokalzeitung.(lesen)

\LDAde{liess}

\item%5 Ich frage mich wirklich, wer ihm dazu (...) (...) (raten; auxiliaire)

\LDAde{geraten; hat}

\item%6 Ich (...) beinahe (...), seinen Bruder einzuladen.(auxiliaire; vergessen)

\LDAde{habe; vergessen}

\item%7 Noch vor ein paar Jahren (...) diese �berfahrt f�r unm�glich (...). (auxiliaire; halten)

\LDAde{w\"urde; gehalten}

\item%8 (...) dir gestern Abend nichts Besonderes (...)? (auxiliaire; auffallen)

\LDAde{W\"urde; aufgefallen}

\item%9 Das (...) sich jetzt leider nicht mehr �ndern. (lassen)

\LDAde{l\"asst}

\item%10 Er (...) in Schwierigkeiten, weil er seine Kr�fte �bersch�tz hatte. (geraten)

\LDAde{ger\"at}

\end{enumerate}

\section*{\LDAex{5}}

\begin{enumerate}
\item%1
  Je pense qu'il a attach\'e beaucoup d'importance \`a cet incident.

\item%2
  M\up{me} Steiner est enrag\'e que les enfants aient fini la barre de chocolat. %TODO Mme

\item%3
Il a d\^u rester ici longtemps; il faut esp\'erer qu'il ne lui soit rien arriv\'e en chemin.

\item%4
  Quand ils ont le temps, ils lisent chaque jour le journal local.
  
\item%5
  J'ai presque oubli\'e d'inviter son fr\`ere.
  
\item%6
Il y a juste quelques ann\'ees, ce trajet en ferry aurait pass\'e pour impossible.  

\item%7
  La soir\'ee d'hier soir ne t'a t'elle pas exceptionnellement marqu\'e.

\item%8
  Je me demande vraiment, qui a pu le lui conseiller.

\item%9
  Malheureusement, cela n'est plus modifiable.
  
\item%10
  Il s'est mis en difficult\'e parce qu'il a surestim\'e sa force.
  
\end{enumerate}

\section*{\LDAtc}

\begin{enumerate}

\item%1 Ne sachant pas lire une carte marine, il s'est �chou� sur un banc de sable dans la baie de Kiel.

  \LDAde{Ohne zu wissen, eine Seekarte zu lesen, hat er sich in der Bucht von Kiel auf eine Sandbank aufgelaufen% [TODO: oder gelaufen?]
.}

\item%2 Il a eu beaucoup de chance que l'homme de barre l'ait vu tomber � l'eau.

\LDAde{Er hat viel Gl\"uck gehabt, dass der Mann am Ruder ihn gesehen hat, als er ins Wasser gest\"urtzt ist % [TODO: oder ihn ins Wasser st\"urtzen gesehen hat?]
.}

\item%3 Le ferry roulait tellement que beaucoup de passagers croyaient qu'ils ne survivraient pas � la travers�e.

\LDAde{Die F\"ahre schlingerte so stark, dass viele Passagiere glaubten, dass sie die \"Uberfahrt nicht \"uberleben w\"urden.}

\item%4 Plusieurs voiliers �taient au mouillage dans la crique quand une temp�te se leva soudain.

\LDAde{Mehrere Segelboote legten vor Anker in der Bucht\LDAverif, wann ein Sturm eingekommen ist.}

\item%5 Quand tu vas en Angleterre, je te conseille de jeter tes pr�jug�s par-dessus bord.

\LDAde{Wenn du nach England f\"ahrst, rate ich dir deine Vorurteilen \"uber Bord zu werfen.}

\item%6 De la chambre d'h�tel o� il passait ses vacances, il pouvait entendre le d�ferlement des vagues.

\LDAde{Von der Zimmer des Gasthof, wo er seine Ferien verbracht, k\"onnte er die Brandung der Welle h\"oren.}
  
\item%7. Son fr�re a fait des sacrifices pendant des ann�es pour s'acheter un voilier.

\LDAde{Sein Br\"uder hat sich einige Jahren geopfert, um ein Segelboot sich einzukaufen.}
    
\item%8 A peine le capitaine avait-il commenc� de donner des ordres pour le sauvetage, que le matelot disparut dans les flots.

\LDAde{Der Kapit\"an h\"atte kaum begonnen, Anweisungen zu geben, um den Matrose zu retten, dass er in der Flugen verschwunden gewesen ist.}

\item%9 Nous lui avions pourtant d�conseill� de vouloir absolument se mesurer � ce gar�on � la t�te en l'air.

\LDAde{Wir h\"atten ihn jenoch davon abgeraten, sich gegen diesen gedankenlosen jungen Mann unbedingt messen zu versuchen.}

\item%10 Lorsque l'�quipe de sauvetage le rep�cha de l'eau glac�e, le naufrag� n'avait presque plus la force d'appeler au secours.

\LDAde{Wenn die Rettungsmannschaft ihn aus eiskaltem Wasser herausfischen hat, h\"atte der Schiffbr\"uchige 
kaum nicht mehr die Kraft um Hilfe zu rufen.} %[Exercice 1]

\item%11 Il faudrait avertir la police, car il serait arriv� depuis bien longtemps si rien ne lui �tait arriv� en route.

  \LDAde{Man soll die Polizei verst\"andigen,
    %[cf Exercice 1(6)]
    denn er schon l\"angst angekommen h\"atte, wenn ihm nichts geschehen h\"atte.}% [Exercice 4(3)]

\item%12 Beaucoup de gens sont d'avis qu'il ne faut pas surestimer l'importance de cet incident.

\LDAde{Viele Leute sind der Meinung, man soll diesen Zwischenfall nicht \"ubersch\"atzen.}% [Exercice 4(1)]

\item%13 Il est furieux parce qu'il ne peut tol�rer qu'on lui dissimule la v�rit�, m\^eme s'il doit en souffrir.

\LDAde{Er is w\"utend, %[Exercice 4(2)]
  weil er dulde es nicht,
  %[Exercice 1(12)]
  dass man vor ihm die Wahrheit verheimlicht, %[1≥2, oder versteckt]
  auch wenn er daran leiden muss.}

\item%14 Si quelque chose de particulier m'avait frappe hier soir, je te l'aurais dit.

\LDAde{Wenn gestern Abend mir etwas Besonderes aufgefallen h\"atte,%[Exercice 4(8)]
h\"atte ich dir davon gesagt.}

\item%15 Il y a quelques annees encore, je n'aurais pas ose rever d'un voyage en Amerique.

\LDAde{Noch vor ein paar Jahren, %[Exercice 4(7)],
  h\"atte ich einmal nicht gewagt, von einer Reise nach Amerika zu tr\"aumen. %[Exercice 1(4)]
}

\end{enumerate}

\section*{\LDAdisk}

\begin{enumerate}

\item%1. Erkl�ren Sie, warum der Schiffbr\"uchige verschwand, wenn ein Fremder auf der Insel auftauchte

\LDAaf

\item%2. Ist die Geschichte "Mann �ber Bord" unmoralisch?

\LDAaf
  
\item%3. Wie k�nnte man Unf�llen voerbeugen, die auf Leichtfertigkeit zur�ckzufhren sind? Geben Sie Beispiele!

\LDAaf
  
\item%4. Sollen Menschen ihr Leben aufs Spiel setzen, um andere Menschen zu retten, die sich leichtfertig in Gefahr begeben haben.

\LDAaf
  
\end{enumerate}

\begin{comment}
\end{comment}

%\chapter{\LDAdemarking{Die Welt der Sicherheit}}

\section*{\LDAex{1}}

%Dans les phrases suivantes, supprimez l'auxiliaire de mode sans modifier le temps du groupe verbal.

\begin{enumerate}

\item %1 Es h\"atte ihm gelungen k\"onnen, wenn er weniger Pech gehabt h\"atte.

\LDAde{Es w\"are ihm gelungen, \ERquoteskip}

\item %2. Er durfte keinen Alkohol mehr trinken 

\LDAde{Er trank keinen Alkohol mehr}

\item %3. Mein altes Klavier hat nicht viel heller klingen k\"onnen als deines

\LDAde{Mein altes Klavier hat nicht viel heller klingen als deines}

\item %4. Wenn die Krise angedauert h\"atte, h\"atten die Roh\"olpreise um wenigstens 6 prozen sinken m\"ussen

\LDAde{Wenn \ERquoteskip, h\"atten die Roh\"olpreise um wenigstens 6\% gesunken}

\item %5. Es sollte ihm endlich gelingen, seine Meinung durchzusetzen

\LDAde{Es war ihm endlich gelingen, \ERquoteskip}

\item %6. Wenn ich meinen Schl\"ussel vergessen h\"atte, h\"atte ich \"uber die Gartenmauer springen m\"ussen.

\LDAde{Wenn \ERquoteskip, h\"atte ich \"uber die Gartenmauer gesprungen}

\item %7 Der Dieb hat in der Menge verschwinden k�nnen

\LDAde{Der Dieb hat in der Menge verschwunden}

\item %8 Niemand konnte ihn zum Nachgeben zwingen

\LDAde{Niemand hat ihn zum Nachgeben gezwungen}

\item %9 Er wollte ins Wasser springen

\LDAde{Er ist ins Wasser gesprungen}

\item %10 Keine andere hat so sch�n singen k�nnen wie sie

\LDAde{Keine andere hat so sh\"on wie sie gesungen}

\item %11 Die K\"alte konnte durch die offen gelassene T\"ur ins Zimmer dringen

\LDAde{Die K\"alte ist durch die offen gelassene T\"ur gedrungen}

\item %12 Als ich j\"unger war, h\"atte ich diesen Roman innerhalb einer Nacht verschlingen k�nnen

\LDAde{Als ich j\"unger war, h\"atte ich diesen Roman innerhalb einer Nacht verschlingen \sidenote{Qui veut dire?!}}

\end{enumerate}

\section*{\LDAex{2}}

%Vous entendez une phrase que vous comprenez mal. Posez une question se rapportant au groupe entre {} pour vous la faire r�p�ter. Vous devez utiliser u ne forme de l'interrogatif "wer"

\begin{enumerate}

\item % 1. Tanja ist mit {einer Schulfreundin} verreist

\LDAde{Mit wem ist Tanja verreist?}

\item % 2. {Wer nicht versucht}, dem kann es nicht gelingen

\LDAde{Wem kann es nicht gelingen?}
 
\item % 3. Ihr Onkel kann an die Ehrlichkeit {der meisten Politiker} nicht glauben

\LDAde{An wessen Ehrlichkeit kann ihr Onkel nicht glauben?}

\item % 4. Sie m�chte nicht an der Stelle {ihrer Mutter} sein

\LDAde{An wessen Stelle m\"ochte sie nicht sein?}

\item %5 Die meisten Autofahrer haben Angst {vor der Polizei}

\LDAaf

\item %6 Er hat sehr lange auf die Antwort {des B\"urgermeisters} gewartet

\LDAaf

\item %7 Ihr Mann wollte {dem Handelsvertreter} nicht glauben

\LDAde{Wem wollte ihr Mann nicht glauben?}

\item %8 Seine Eltern hatten {mit einem zus\"atzlichen Gast} nicht gerechnet

\LDAde{Mit wem hatten seine Eltern nicht gerechnet?}

\end{enumerate}

\section*{\LDAex{3}}

%Dans les phrases suivantes, remplacer le sujet par une forme du relatif "wer" pour g�n�raliser.  
% Attention: 
%  a. La g�n�ralisation implique parfois un passage au temps pr�sent
%  b. Dans certains cas, le relatif doit �tre repris dans la principale par "der"
			
\begin{enumerate}

\item % 1. Manfred, der viel Geld hat, kann sich alles leisten

\LDAde{Wer viel Geld hat, \ERquoteskip}

\item % 2. Herr M\"uller, dem die Papiere gestohlen wurden, konnte nicht aus der DDR herauskommen

\LDAde{Wer seine Papiere gestohlen wurden, \ERquoteskip}

\item % 3. Sein Gro�vater, der fr\"uher genug gespart hatte, brauchte sich keine Sorge zu machen

\LDAde{Wer fr\"uher gespart hatte, brauchte sich keine Sorgen zu machen}

\item % 4. Reiner, dem es gelungen ist, beim Lotto die 6 richtigen Zahlen anzukreuzen, hat mehr gewonnen, als er je verdienen k�nnte

\LDAde{Wem ist es gelungen, beim Lotto die 6 richtigen Zahlen anzukreuzen, und so mehr gewonnen hat, als er je verdienen k\"onnte?}

\item %5 Walther, der Neapel noch nicht gesehen hat, hat wenigstens die M�glichkeit, davon zu tra\"umen

\LDAde{Wer hat noch nicht Neapel gesehen, und wenigstens die M\"oglichkeit hat, davon zu tra\"umen?}

\item %6 J\"urgen, der weder sparen noch warten kann, kauft vieles auf Kredit

\LDAde{Wer kann weder sparen noch warten, und deshalb kauf vieles auf Kredit?}

\item %7 Ihr Sohn, dem niemand vertrauen will, glaubt selber nicht mehr an sich

\LDAde{Wem will niemand vertrauen, und glaubt selber nicht mehr an sich?}

\end{enumerate}

\section*{\LDAex{4}}

%Compl�tez les phrases suivantes par "bis" en y associant, si n�cessaire, la pr�position ou ou l'adverbe qui convient. 

\begin{enumerate}

\item % 1. Er hat dieses Haus ... seinem Tod bewohnt

\LDAde{bis zu}

\item % 2. Ein deutscher Autofahrer hat uns ... spanischen Grenze mitgenommen

\LDAde{bis an die}

\item % 3. Alle Teilnehmer ... einen waren mit dem Vorschlag einverstanden

\LDAde{bis auf}

\item % 4. Du kannst noch lange warten, ... sie Zeit hat, darum zu k\"ummern

\LDAde{bis}

\item % 5. Ihr Sohn hat es nicht einmal ... Abitur gebracht

\LDAde{bis zum}

\item % 6. Sie d\"urfen ihre Freunde nicht ... Grenze begleiten

\LDAde{bis an die}

\end{enumerate}

\section*{\LDAex{5}}

% Traduisez les phrases que vous avez faites dans l'exercice pr�c�dent.

\begin{enumerate}

\item % 1. Er hat dieses Haus ... seinem Tod bewohnt

Il a habit� dans cette maison jusqu'� sa mort.

\item % 2. Ein deutscher Autofahrer hat uns ... spanischen Grenze mitgenommen

Un automobiliste allemand nous a emmen� jusqu'� la fronti�re espagnole.

\item % 3. Alle Teilnehmer ... einen waren mit dem Vorschlag einverstanden

Tous les participants, sauf un, �taient d'accord avec la proposition.

\item % 4. Du kannst noch lange warten, ... sie Zeit hat, darum zu k\"ummern

Tu peux encore attendre, avant qu'elle n'ait le temps de s'en soucier.

\item % 5. Ihr Sohn hat es nicht einmal ... Abitur gebracht

Son fils (� elle) n'a pas atteint le niveau du bac.

\item % 6. Sie d\"urfen ihre Freunde nicht ... Grenze begleiten

Vous ne devez pas accompagner votre ami jusqu'� la fronti�re.

\end{enumerate}

\section*{\LDAex{6}}

% Traduisez les phrases suivantes en employant "immer" associ� � un comparatif

\begin{enumerate}

\item % 1. Il fait de plus en plus froid

\LDAde{Es wird immer k\"alter.}

\item % 2. Il d�pense de plus en plus d'argent

\LDAde{Er gibt immer mehr Geld aus.}

\item % 3. Je crois de moins en moins � la guerre nucl�aire

\LDAde{Ich glaube immer wenigen an den Atomkrieg.}

\item % 4. De moins en moins de gens ont un id�al de vie.

\LDAde{Immer weniger Leute haben ein Lebensideal.}

\item % 5. La qualit� s'am�liore de plus en plus.

\LDAde{Die Qualit\"at wird immer besser.}

\item % 6. La vie de devient de plus en plus rude.

\LDAde{Das Leben wird immer mehr schwierig.}

\item % Dieter devient de plus en plus malhonn�te.

\LDAde{Dieter wird immer mehr unehrlich.}

\item % Il conduit avec de plus en plus de prudence.

\LDAde{Er f\"ahrt mit immer mehr Vosicht.}

\item % Il criait de plus en plus fort

\LDAde{Er schreie immer la\"uter.}

\item % Il boit de plus en plus

\LDAde{Er trinkt immer mehr.}

\item % Les prix baissent de plus en plus vite

\LDAde{Die Preisen abnehmen immer schneller.}

\item % Les Allemands vivent de plus en plus sainement

\LDAde{Die Deutschen f\"uhren immer mehr ein gesundes Leben.}

\item % Il vient nous voir de plus en plus rarement

\LDAde{Er besucht uns immer selten.}

\item % Je trouve le temps de plus en plus long

\LDAde{Die Zeit wird mir immer l\"anger.}

\item % Cela me para\^it devenir de plus en plus difficile

\LDAde{Es wird mir immer schwieriger.}

\item % J'y crois de moins en moins

\LDAde{Das glaube ich immer weniger.}

\item % Je trouve son histoire de plus en plus ennuyeuse

\LDAde{Seine Geschichte wird mir immer langweilig.}

\end{enumerate}

\section*{\LDAtc}


\begin{enumerate}

\item %1. S'il avait put pr�voir ce qui est arriv�, il ne serait certainement pas venu

\LDAde{H\"atte er vorgesehen was ihm geschehen is, h\"atte er wirklich nicht gekommen.}

\item %2. Je ne sais toujours pas sur quoi il fonde son espoir d'un avenir meilleur

\LDAde{Ich weiss noch nicht, \LDAaf.}

\item %3. Stefan Weig a grandi � une �poque o� tout semblait �tre fond� pour durer

\LDAde{Stefan Weig ist an einer zeit aufgewachsen, in dem alle gegr\"undet zur dauern schien.}

\item %4. � cette �poque l�, qu'il d�crit comme l'�ge d'or de la s�curit�, il y avait en Europe une stabilit� dans tous les domaines de la vie sociale et politique

\LDAde{Zu dieser Zeit, die er als ein goldenes Zeitalter der Sicherheit beschreibt, gibt es Stabilit\"at in allen Bereiche des sozialen und politischen Lebens.}

\item %5. Pour jouir pleinement du pr�sent, il ne suffit pas de pouvoir regarder l'avenir en toute qui�tude, car des bouleversements sont toujours possibles qui d�truisent brutalement les bases de la soci�t�

\LDAde{Um die Gegenwart v�llig zu genie�en, ist es nicht genug, die Zukunft sorglos blicken zu k�nnen, denn radikale Krisen entstehen k�nnen, die die Gr\"undungen der Geselschaft st\"urzen.}

\item %6. � sa place, je ne la croirais pas, car elle a d�j� plusieurs fois abus� de sa confiance

\LDAde{An seiner Stelle, w\"urde ich ihr nicht glauben, als sie seiner Vertrauen mehrmals missbraucht hat.}

\item %7. Il me para�t difficile de ramener leurs points de vue � un d�nominateur commun, car ils n'ont vraiment rien en commun.

\LDAde{Es scheint schwierig, ihre Standpunkte aug einer gemeinsamer Nennen zu bringen} % Lexique 10/10

\item %8. Beaucoup de gens triment toute leur vie pour faire des �conomies qui sont englouties d�s que les prix augmentent.

\LDAde{Viele Leute arbeiten das ganzen Leben um Erspanisse% Haben oder Sein, 1?2 
zu machen, die verschwunden werden, sobald die Preise zunehmen werden}

\item %9. Les Fran�ais sont en majorit� convaincus de son honn�tet� et croient en sa volont� de changer de politique.

\LDAde{Die Franzosen sind von seiner Ehrlichkeit \"uberzeugt, und glauben an seinen Wille, die Politik zu \"andern.}

\item %10. Il a promis de lib�raliser la vie politique du pays s'il r�unissait � arriver au pouvoir.

\LDAde{Er hat versprochen das Politischen Leben des Landes zu offen, ob er es geschafft w\"urde,  an die Macht anzukommen} % Der Spielverderber | Die Gesellschaft im Wandel

\item %11. Il ne faut pas trop croire aux promesses des hommes politiques; ce ne sont souvent que de belles paroles creuses.

\LDAde{Man soll die Versprechen der Politiker nicht zu viel glauben; die sind oft sch\"one leeren Worte}

\item %12. Quiconque a lu les romans de Zola ne peut oublier de si t�t la pauvret� du peuple au XIX� si�cle.

\LDAde{Wer die von Zola gelesen hat, kann nicht zu fr\"uh vergessen, die Armut des Volkes des 19. Jahrhunderts vergessen}

\item %13. Pour mener une politique de gauche, le gouvernement a d�cid� de nationaliser les banques et quelques entreprises importantes pour l'�conomie du pays.

\LDAde{Um eine Linkspolikit zu f\"uhren, hat die Regierung entscheiden, die Banken und einige gewaltigen Firmen zu nationalisieren}

\item %14. Bien qu'il soit au pouvoir depuis plusieurs ann�es, ce parti n'a pas r�ussi � promouvoir la politique progressiste dont r�vaient ses adh�rents.

\LDAde{Obwohl er vor mehrere Jahren an die Macht  gestanden ist, hat dieser Partei nicht geschafft, den links Politischen Kurs zu f\"ordern, von dem seine Anh\"anger tra\"umten}

\end{enumerate}

\section*{\LDAdisk}

\begin{enumerate}

\item %1. Der Mensch und seine Vergangenheit

\LDAaf

\item %2. Kann man im politischen Bereich von Fortschritt sprechen?

\LDAaf

\item %3. Memoiren von Schriftstellern und Geschichtsunterricht

\LDAaf

\item %4. Kann die heutige Armut in der Dritten Welt mit der des 19. Jahrhunderts in Europa verglichen werden?

\LDAaf

\item %5. Sind wir innerhalb Europas vor R�ckf�llen in die Barbarei sicher?

\LDAaf

\end{enumerate}

%---
%Bonus
%---

% Das Recht(, e) und die Pflicht(, en) / Le droit et le devoir
% Der DM Umlauf / la circulation du DM
% Verb�rgte damit ihre Unwandelbarkeit / garantissait ainsi leur inalt�rabilit�
% Jeder wu�te, wieviel er besa� oder wieviel ihm zukam / chacun savait combien il poss�dait ou lui �tait redevable
% Ma� und Gewicht / Poids et mesure
% Jede Familie hatte ihr bestimmte Budget bereitgestellt / Chaque famille mettait de c�t� un budget pr�cis
% Hof und Gesch�ft vererbten sich von Geschlecht zu Geschlecht / Fermes et magasins se transmettaient de g�n�ration en g�n�ration
% w�hrend ein S�ugling in der Wiege lag / pendant qu'un enfant avait � l'�ge du berceau
% schien bereits unm�glich in einem Zeitalter der Vernunft / semblait m�me impensable dans une �re de la raison
% Dieses Gef�hl der Sicherheit war der Besitz von Millionen / ce sentiment de s�curit� �tait la possession (privil�ge) de millions
% weitere Kreise begehrten ihren Teil an diesem kostbaren / d'autres cercles (groupes sociaux) r�clamaient leur part de ce bien pr�cieux
% die Besitzenden / les poss�dants
% k�nftige Generationen / les futures 
% Nur wer sorglos in die Zukunft blicken konnte, geno� mit gutem Gef�hl die Gegentwar / seul 
% celui qui pouvait se projeter 
% dans l'avenir avec confiance, profite du pr�sent avec s�r�nit�
% entstehen order bestehen / se produire ou �tre en place

%\chapter{Leben mit Chemie}

\begin{Exercise}[title={\nameref{O-chemie:hyp}}, label={chemie:hyp}]
Continuez la phrase en int\'egrant les \'el\'ements entre parenth�ses dans une phrase hypoth\'etique. 

\Question Wenn ich an seiner Stelle gewesen w�re, \ldaBLANK \\(darauf / keine / nehmen / ich / R�cksicht)

\ERquoteskip, h\"atte ich darauf keine R\"ucksicht genommen.

\Question Ohne die Erbschaft seines Onkel vor zehn Jahren, \ldaBLANK\\(ein / arm / bleiben / Teufel / er)

\ERquoteskip, w\"are er ein armes Teufel geblieben.

\Question Wenn er geschickter gewesen w�re, \ldaBLANK\\ 
(seine / er / Forderung / durchsetzen / k�nnen )

\ERquoteskip, h\"atte er seine Forderung durchsetzen k\"onnen.

\Question Wenn Herr M�ller mehr Erfahrung gehabt h�tte, \ldaBLANK\\(er / scheitern / miserabel / in / sein / nicht / Beruf / so)

\ERquoteskip, h\"atte er nicht so miserabel in seinem Beruf gescheitert 

\item Er h�tte die richtige L�sung finden k�nnen, \ldaBLANK\\
(von  / an / er / Anfang / nicht / sich / irren )

\ERquoteskip, wenn von Anfang an nicht geirrt h\"atte

\Question Helmut w�re ein kleiner Bankangestellter geblieben, \ldaBLANK
(ihm / sein / wenn / Schwiegervater / diese / Stelle / zu / verhelfen / nicht )

\ERquoteskip, wenn sein Schwiegervater ihm zu diese Stelle nicht verholfen h\"atte.

\Question Wenn du wolltest, \\(ich / begleiten / dich / gern)

\ERquoteskip, w\"urde ich dich gern begleiten.

\item Um seinen Plan zu verwirklichen, \ldaBLANK\\(er / m�ssen / auf / die / moderne Gesellschaft / die Vorteile / verzichten)

\ERquoteskip, muss er auf die Vorteile der moderne Geselschaft verzichten 

\item Vom Standpunkt der Gr�nen aus, \ldaBLANK\\
(sein / ein / aus / die Kernenergie / Ausstieg / m�glich / durchaus)

\ERquoteskip, ist ein Ausstieg aus der Kernenergie durchaus m\"oglich

\Question Wenn die Menschen nicht so kurzsichtig w�ren, \ldaBLANK
(der Umweltschutz / der Vorrang / geben / werden)

\ERquoteskip, w\"urde den Umweltschutz den Vorrang geben

\end{Question}

\begin{Answer}[ref={O-sein:gelingen}]
\end{Answer}

% Traduisez les phrases que vous avez faites dans l'exercice pr\'ec\'edent.
\begin{comment}

\section*{\LDAex{2}}

% Traduisez les phrases que vous avez faites dans l'exercice pr\'ec\'edent.

\begin{enumerate}

\item Si j'\'etais � sa place, je n'y accorderais aucune consid\'eration

\item Sans l'h\'eritage de son oncle il y a 10 ans, il serait rest\'e un pauvre diable.

\item S'il avait \'et\'e plus adroit, sa revendication aurait pu aboutir

\item Si M. M\"uller avait eu plus d'exp\'erience, il n'aurait pas \'echou\'e si mis\'erablement dans son travail.

\item Il aurait trouv\'e la bonne solution, s'il ne s'\'etait pas tromp\'e d\`es le d\'ebut.

\item Helmut serait rest\'e un petit employ\'e de banque, si son beau p\`ere ne l'vait pas aid\'e � obtenir cette place.

\item Si tu le voulais, je t'accompagnerais volontiers.

\item Afin de r\'ealiser son projet, il doit abandonner les avantages de la vie moderne.

\item Du point de vue des verts, l'abandon de l'\'energie nucl\'eaire est tout � fait possible.

\item Si les gens n'\'etaient pas si born\'es, ils donneraient la priorit\'e � la protection de l'environnement.

\end{enumerate}

\section*{\LDAex{3}}

% Compl\'etez les phrases suivantes par une forme de wissen, k�nnen ou kennen. 

\begin{enumerate}
\item % 1. Wer (...), wie lange es noch dauern wird, bis er zur�ckkommen (...)

weiss / k\"onnte

\item % 2. Ich (...), da� er dir geholfen h�tte, wenn er (...) h�tte, denn ich (...) ihn sehr gut.

weiss / k\"onnen / kenne

\item % 3. Ich (...) mir nicht vorstellen, da� er das schon (...), denn er hat keinen Kommentar gegeben.

kann / weiss 

\item % 4. Ich wundere mich kaum �ber sein dreistes Verhalten, denn schon als Kind (...) er keine Grenzen.

wusste

\item % 5. Wer (...) schon sagen, ob die Vorschriften streng genug sind, um einen Umfall zu verhindern.

kann

\item % 6. Nicht, da� ich (...)!

weiss

\item %7. Jeder Mensch (...) die kleine Nacthmusik von Mozart

kennt

\item % 8. Ich (...) nicht (...), da� er davon profitieren w�rde

konnte / wissen

\item % 9. Wenn du glaubst, da� ich das nicht (...), da (...) du mich aber schlecht!

kann / kennst

\item % 10. Ich habe schon immer (...), da� es so kommen w�rde

gewusst

\item % 11. Ich habe nicht fr�her kommen (...)

k\"onnte

\item % 12. (...) du, ob sie ihn pers�nlich (...) oder nur vom Sehen

Weisst / kennt

\end{enumerate}

\section*{\LDAex{4}}

\begin{enumerate}

\item % 1. 

Qui sait combien de temps peut encore s'\'ecouler, avant qu'il ne puisse revenir.

\item % 2. 

Je sais qu'il t'aurais aid\'e, s'il avait pu, car je le connais tr\`es bien.

\item %3

Je ne peux m'imaginer qu'il le sache d\'ej�, parce qu'il n'a pas fait de commentaire

\item %4. 

Je ne m'en ferais � peine pour l'imprudence de son comportement, car d\'ej� enfant, il ne connaissait pas de limite.

\item %5

Qui peut dire, si les r\`egles sont assez s\'ev\`eres pour emp\^echer un accident?

\item % 6. 

Non, pas que je sache!

\item % 7. 

Chaque homme connait la petite musique de nuit de Mozart

\item % 8. 

Je ne pouvais pas savoir qu'il en profiterait

\item % 9. 

Que tu crois que je n'en sois pas capable, tu me connais mal!

\item % 10. 

Je savais d\'ej� qu'il viendrait ainsi

\item % 11. 

Je n'ai pas pu venir plus t�t

\item % 12. 

Sais tu si elle le conna�t personnellement, ou seulement de vue?

\end{enumerate}

\section*{\LDAex{5}}

%Traduisez les phrases suivantes.

\begin{enumerate}

\item %1. Vouloir, c'est pouvoir.

Wollen is macht.

\item %2. D'o� sais-tu qu'il ne sait pas nager?

Woher weisst du, dass er nicht schwimmen kann?

\item %3. Je ne savais pas qu'il viendrait aussi.

Ich wusste nicht, dass er auch kommen w\"urde.

\item %4. Sa femme sait tr�s bien s'affirmer

Seine Frau weiss ganz schon wie sich zu behaupten.

\item %5. Si j'avais su cela, je ne serais pas venu

Wenn ich es wusste, w\"are ich nicht gekommen.

\item %6. Je ne pouvais pas savoir qu'il renoncerait si vite

Ich k\"onnte nicht wissen, dass er so schnell verzichten w\"urde 

\item %7. Autant que je sache, aucune mesure n'a \'et\'e prise

Soweit ich weiss, keine Massnahme getroffen ist

\item %8. Savez-vous si la production actuelle d'\'electricit\'e suffira dans cinq ans?

Wissen Sie, ob die gegenw\"artige Stromproduktion in f\"unf Jahren ausreichen wird?

\item %9. Je ne le connais que de vue.

Ich kenne ihn nur vom Sehen.

\item %10. Il ne sait pas ce qu'il dit.

Er weiss nicht, was er sagt.

\item %11. J'ai tr�s bien connu votre p�re.

Ich habe ihren Vater kennengelernt.

\item %12. Savez-vous, si cela est (serait) possible?

Wissen Sie ob es m\"oglich ist (w\"are)?

\item %13. Je ne connais que la vieille ville.

Ich kenne nur die Altstadt.

\item %14. Si seulement il pouvait oublier.

Wenn er nur vergessen k\"onnte!

\item %15. Personne ne sait ce qu'il est devenu.

Niemand weiss was aus ihm geworden ist

\item %16. Je ne connais cette possibilit\'e technique que par ou�-dire mais je sais qu'elle existe.

Ich kenne diese technische M\"oglichkeit nur vom H\"orensagen, aber ich weiss, dass es existiert

\item %17. Je saurai bien le trouver.

Ich werden wissen ihn zu finden.

\end{enumerate}

\section*{\LDAtc}

\begin{enumerate}

\item % 1. M�me les directives et les mesures les plus s\'ev�re sont peine perdue si les gens ne sont pas pr�ts � les respecter---comme ces \'eleveurs de Bretagne qui mettaient des hormones interdites dans les aliments.

Auch strenge Verschriften und Massnahmen w\"urden hilflos sein, wenn die Leute denen nicht befolgen---als diesen Viehz\"uchter in der Bretagne verbotene Hormone in dem Futtermittel einsetzen

\item % 2. En RFA, apr�s l'accident de Tchernobyl, beaucoup d'hommes politiques ont, contrairement � la France, r\'eclam\'e un abandon de l'\'energie nucl\'eaire.

In der BRD, nach dem Unfall von Tchernobyl, viele Politiker haben, anderseits Frankreich, einen Ausstieg aus der Kernenergie gefordet. 

\item % 3. Le premier ministre voudrait bien donner la priorit\'e � la protection de l'environnement, mais il ne le font pas par peur des cons\'equences politiques d'une telle d\'ecision.

Der Ministerpr\"asident w\"urde gern dem Umweltschutz das Vorrang geben, aber das tun sie nicht vor Angst politischen Folgen von einer solcher Entscheidung.

\item % 4. Depuis que nous nous sommes perdus de vue, je ne sais absolument pas ce qu'il est devenu

Seitdem wir uns aus den Augen verloren haben, weisse ich nicht was aus ihm geworden ist.

\item % 5. Comme nous ne pouvons pas renoncer � la croissance \'economique sans compromettre l'emploi dans l'industrie, la seule issue consiste en une industrie diff\'erente.

\item % 6. Chacun sait que l'utilisation de la chimie dans l'agriculture provoque des effets secondaires d\'esastreux, en particulier pour la nappe phr\'eatique.

\item % 7. Presque tous les produits agricoles que nous consommons quotidiennement contiennent des d\'echets chimiques r\'esiduels plus ou moins toxiques.

\item % 8. Par manque de confiance en ses collaborateurs, il pr\'ef�re tout faire lui m�me.

\item % 9. Si les hommes n'avaient pas la vue aussi courte, nous aurions certainement moins de confort, mais notre qualit\'e de vie ne serait pas n\'ecessairement plus mauvaise.

\item % 10. J'ai eu beaucoup de change que le p�re d'un ami m'ait fait obtenir cet emploi dans une entreprise allemande dans laquelle l'aspect humain de chaque probl�me est pris prioritairement en consid\'eration.

\item % 11. Je ne sais pas ce que j'aurais fait si j'avais \'et\'e � sa place, mais je n'aurais pas non plus renonc\'e � mon projet.

\item % 12. Il a essay\'e de r\'ealiser seul son projet, mais il a \'echou\'e lamentablement par manque d'exp\'erience

\end{enumerate}

\section*{\LDAdisk}

\begin{enumerate}
\item %1. Wachstum um jeden Preis?

\item % 2. Die Januk�pfigkeit der Inderstialisierung in der Dritten Welt.

\item % 3. Wohlstand und Lebensqualit�t

\item % 4. Die Priorit�ren der heutigen Industriegesellschaft

\item % 5. Eine saubere Umwelt: reine Utopie oder eine Frage der Zeit?

\item % 6. Die industrielle Herausforderung

\item % 7. Ist der Mensch seiner Aufgabe in der modernen Industriegesellschaft gewachsen?

\end{enumerate}

\end{comment}

%Bonus
% Auf etwas nutzen nehmen / tirer profit de qq chose
% Auf die Risiken hinweisen
% Eine Herausforderung gewachsen sein

\chapter{\nameref{O-ch:seinoder}}

\begin{Exercise}[title={\nameref{O-sein:gelingen}}, label={sein:gelingen}]
Traduisez les phrases suivantes:
\Question Si seulement je pouvais r\'eussir \`a l'en convaincre!
\Question Tu ne r\'eussiras jamais \`a faire assez d\'economies pour t'acheter cette voiture de luxe
\Question Comment veux tu qu'il r\'eussisse \`a r\'ealiser du profit?
\Question Le nouveau mod\`ele est assez r\'eussi
\Question Leur m\`ere a enfin r\'eussi \`a ne plus se laisser manipuler par la publicit\'e
\Question Elle n'a malheureusement pas r\'eussi sa surprise
\Question Celui qui se laisse manipuler ne peut pas r\'eussier \`a renoncer \`a consommer
\Question Je ne suis pas arriv\'e \`a mettre de l'ordre dans ses papiers
\end{Exercise}

\begin{Answer}[ref={O-sein:gelingen}]

\Question % 1. Si seulement je pouvais r\'eussir \`a l'en convaincre!

Wenn es mir nur gelungen w\"are, ihm davon zu \"uberzeugen

\Question % 2. Tu ne r\'eussiras jamais \`a faire assez d'\'economies pour t'acheter cette voiture de luxe

Es wird dir niemals gelingen, genug zu sparen, um dir dieses Luxusauto zu kaufen

\Question % 3. Comment veux-tu qu'il r\'eussisse \`a r\'ealiser du profit?

Wie stellst du dir vor, dass es ihm gelingen w\"urde, einen Gewinn zu machen.

\Question % 4. Le nouveau mod\`ele est assez r\'eussi

Das neue Modelle is erfolgreich

\Question % 5. Leur m\`ere a enfin r\'eussi \`a ne plus se laisser manipuler par la publicit\'e

Es ist ihrer Mutter endlich gelungen,
sich nicht mehr durch Werbung manipuliert zu werden

\Question % 6. Elle n'a malheureusement pas r\'eussi sa surprise

Leider ist ihr \"uberraschung nicht gelungen

\Question % 7. Celui qui se laisse manipuler ne peut pas r\'eussir \`a renoncer \`a consommer

Wem sich manipuliert l\"asst, kann es nicht gelungen, konsumieren zu verzichten

\Question % 8. Je ne suis pas arriv\'e \`a mettre de l'ordre dans ses papiers

Es ist mir nicht gelungen, Ordnung in meiner Papiere zu setzen.

\end{Answer}


\begin{Exercise}[title={\nameref{O-sein:aussen}},label={sein:aussen}]

Compl\'etez les phrases suivantes par "aussen", "\"ausserlich", "draussen", "ausser", "\"ausser...", "ausserhalb, "innen", "innerlich", "drinnen", "inner..." ou "innerhalb".

\Question Dem \ldaBLANK minister steht die Polizei zur Verf\"ugung \ldaBLANK Ordnung zu gew\"ahrleisten

\Question \ldaBLANK ist das Haus nicht ohne \"ahnlichkeit mit einem Bunker, aber man vergisst es schnell, wenn man \ldaBLANK sitzt, denn es ist sehr gem\"utlich

\Question Seine Frau will \ldaBLANK der Stadt wohnen, damit kie Kinder einen Garten zum Spielen haben

\Question Er hat mit versprochen, nicht eigensinnig zu sein und sich den \ldaBLANK Umst\"anden anzupassen

\Question Kannst du so lange \ldaBLANK warten ?

\Question Da er als \ldaBLANK minister f\"ur die \ldaBLANK Angelegenheiten nat\"urlich nicht zust\"andig ist, hat er sich geweigert, auf eine Frage \"uber die Jugendkriminalit\"at zu antworten

\Question Die \ldaBLANK stadt ist versuchsweise f\"ur den Autoverkehr gesperrt worden.

\Question \ldaBLANK ihrem Beruf haben viele Menschen keine anderen Interessen im Leben

\Question Die Gesetze gelten \"uberall \ldaBLANK der Landesgrenzen

\Question Der Artz empf\"angt keine Patienten \ldaBLANK der Sprechstundent

\end{Exercise}
\begin{comment}

\begin{Answer}[ref={O-sein:aussen}]

% Compl\'etez les phrases suivantes par "aussen", "\"ausserlich", "draussen", "ausser", "\"ausser...", "ausserhalb, "innen", "innerlich", "drinnen", "inner..." ou "innerhalb".

\Question aussen, innere

\Question Von aussen, drinnen

\Question ausserhalb

\Question aussen

\Question draussen

\Question aussen, innerlichen

\Question innere

\Question ausser

\Question innerhalb

\Question ausserhalb

\end{Answer}

\begin{Exercise}[label={sein:aussen2}]

Traduisez les phrases de l'exercice pr\'ec\'edent.

\end{Exercise}

\begin{Answer}[ref={O-sein:aussen2}]

\Question % 1. 

La police est \`a la disposition de la police pour assurer l'ordre

\Question % 2. 

De l'ext\'erieur, cette maison n'est pas sans ressemblance avec un bunker, mais on l'oublie, une fois \`a l'int\'erieur, qui est tr\`es confortable

\Question % 3. 

Sa femme veut habiter en dehors de la ville, pour que ses enfants aient un jardin dans lequel jouer

\Question % 4.

Il m'a promis de ne pas \^etre obstin\'e et de s'adapter aux circonstances ext\'erieures

\Question % 5. 

Peux tu attendre si longtemps dehors?

\Question % 6. 

Parce qu'il est ministre de l'ext\'erieur, il n'est pas responsable des affaires int\'erieures, il a refus\'e de r\'epondre \`a une question sur la criminalit\'e des jeunes.

\Question % 7. 

Le centre ville est temporairement ferm\'e \`a la circulation automobile.

\Question % 8.

En dehors de leur profession, beaucoup de gens n'ont pas d'autre int\'er\^et dans la vie

\Question % 9. 

Les lois valent \`a l'int\'erieur des fronti\`eres

\Question % 10. 

Le m\'edecin ne re�oit pas de patient en dehors des heures annonc\'ees

\end{Answer}

\begin{Exercise}[title={\nameref{O-sein:was}}, label={sein:was}]

Compl\'etez les phrases suivantes par "was" et/ou un ant\'ec\'edent logique. Attention: vous devez mettre un mot par blanc

\Question 1. \ldaBLANK er getan hat, \"uberrascht mich kaum
\Question 2.  Das ist wirklich \ldaBLANK \ldaBLANK, \ldaBLANK ich dazu sagen kann
\Question 3. Du kannst froh sein, \ldaBLANK zu haben, \ldaBLANK du brauchst 
\Question 4. Seine Papiere zu verlieren ist \ldaBLANK \ldaBLANK, \ldaBLANK einem passieren kann, wenn man im Ausland ist
\Question 5. Das ist \ldaBLANK, \ldaBLANK ich zur Zeit brauche
\Question 6. Es gibt nicht \ldaBLANK, \ldaBLANK man ihm vorwerfen kann
\Question 7. Es ist nicht \ldaBLANK Gold, \ldaBLANK gl\"anzt

\end{Exercise}
\begin{Answer}[ref={O-sein:was}]

\Question % 1. \ldaBLANK er getan hat, \"uberrascht mich kaum
Was
\Question % 2.  Das ist wirklich \ldaBLANK \ldaBLANK, \ldaBLANK ich dazu sagen kann
nicht, alles, was
\Question % 3. Du kannst froh sein, \ldaBLANK zu haben, \ldaBLANK du brauchst 
alles, was
\Question % 4. Seine Papiere zu verlieren ist \ldaBLANK \ldaBLANK, \ldaBLANK einem passieren kann, wenn man im Ausland ist
leider, wirklich, was
\Question % 5. Das ist \ldaBLANK, \ldaBLANK ich zur Zeit brauche
alls, was
\Question % 6. Es gibt nicht \ldaBLANK, \ldaBLANK man ihm vorwerfen kann
viel, das
\Question % 7. Es ist nicht \ldaBLANK Gold, \ldaBLANK gl\"anzt
alles, was

\end{Answer}

\begin{Answer}[ref={spiel:theme}]

\begin{Exercise}[title={\ldaTHEME}, label={sein:theme}]

\Question Si les gens n'achetaient que ce dont ils ont vraiment besoin, il y aurait beaucoup moins de gaspillage

\Question Quand le march\'e est satur\'e, la production exc\`ede tr\`es vite les besoins

\Question Pour lancer ce nouveau produit qui ne correspond pas \`a un besoin, cette entreprise a d� s'adresser \`a une agence publicitaire

\Question Dans le prix de vente de beaucoup de produits, la part des frais de publicit\'e et de l'emballage est souvent plus \'elev\'e que celle de la fabrication proprement dite

\Question Mon oncle d'Am\'erique m'a assur\'e qu'il changeait tous les ans de voiture pour ne pas se faire remarquer

\Question Le consommateur moyen ne se doute pas que ses besoins sont cr\'e\'es par la publicit\'e sans laquelle la soci\'et\'e moderne est inconcevable

\Question Au lieu de faire r\'eparer son viel appareil de t\'el\'evision, il pr\'ef\`ere en acheter un neuf \`a cr\'edit

\Question Sachant qu'il a tout \`a profusion chez lui, je ne sais pas quoi lui offrir \`a son anniversaire

\Question En p\'eriode de crise, les agens de publicit\'e se frottent les mains

\Question Ses coll\`egues de travail l'\'evitent, car il s'est rendu suspect en venant sa voiture pour limiter ses d\'epenses

\Question Le pauvre: apr\`es avoir \'et\'e expuls\'e de sons propre pays parce qu'il vantait sans cesse les avantages de la soci\'et\'e de consommation occidentale, il vit maintenant de l'aide sociale en RFA et il ne peut rien acheter hormis les choses d'usage courant.

\Question Je ne pourrai jamais me pardonner de ne m'�tre pa tu

\Question Apr\`es avoir vendu son magasin de fournitures de bureau, sa femme a ouvert une boutique de mode

\Question Il est soup�onn\'e d'appartenir \`a une organisation terroriste parce qu'il a occup\'e une maison vide avec quelques autres \'etudiants en sociologie.

\end{Exercise}


\begin{comment}

\begin{Answer}[ref={O-sein:theme}]

\Question % 1. Si les gens n'achetaient que ce dont ils ont vraiment besoin, il y aurait beaucoup moins de gaspillage

Wenn die Leute kauften was sie nur wirklich brauchten, w\"urde viel weniger verschwendet.

\Question % 2. Quand le march\'e est satur\'e, la production exc\`ede tr\`es vite les besoins

Wenn der Markt nicht ges\"attigt ist, \"ubersteigt die Produktion sehr schnell den Bedarf

\Question % 3. Pour lancer ce nouveau produit qui ne correspond pas \`a un besoin, cette entreprise a d� s'adresser \`a une agence publicitaire

Un dieses neue Produkt auf den Markt zu bringen, d\"urfte diese Firma eine Werbeagentur beauftragen.

\Question % 4. Dans le prix de vente de beaucoup de produits, la part des frais de publicit\'e et de l'emballage est souvent plus \'elev\'e que celle de la fabrication proprement dite

In der Preis vieler G\"uter, ist der Teil der Werbung und Verpackung oft gr\"osser als derjenige der Herstellung.

\Question % 5. Mon oncle d'Am\'erique m'a assur\'e qu'il changeait tous les ans de voiture pour ne pas se faire remarquer

Mein Unkel aus Amerika hat mir versichert, dass er jedes Jahres sein Auto w\"achselt, um nicht auffallen zu werden% [Text: Wer in Amerika keinen Fernsehapparat besitzt, f\"allt auf]

\Question % 6. Le consommateur moyen ne se doute pas que ses besoins sont cr\'e\'es par la publicit\'e sans laquelle la soci\'et\'e moderne est inconcevable

Der Konsument verzweifelt sich nicht, dass seine Bed\"urfnisse von der Werbung geweckt sind, ohne derer ist die moderne Gesellschaft undenkbar.  

\Question % 7. Au lieu de faire r\'eparer son viel appareil de t\'el\'evision, il pr\'ef\`ere en acheter un neuf \`a cr\'edit

Anstatt seinen alten Fernsehapparat zu reparieren, will er lieber einen neuen kaufen.

\Question % 8. Sachant qu'il a tout \`a profusion chez lui, je ne sais pas quoi lui offrir \`a son anniversaire

Wissend, dass er zu Hause alles in H\"ulle un F\"ulle hat, ich weiss nicht, was ihm zu seiner Geburtstag zu anbieten.

\Question % 9. En p\'eriode de crise, les agens de publicit\'e se frottent les mains

In Zeiten der Krise, reiben sie sich die H\"ande

\Question % 10. Ses coll\`egues de travail l'\'evitent, car il s'est rendu suspect en venant sa voiture pour limiter ses d\'epenses

Seine Berufskollegen meiden ihn, weil er sich verd\"achtig, indem er sein Auto verkauft hat, um seine Ausgabe abzunehmen 

\Question % 11. Le pauvre: apr\`es avoir \'et\'e expuls\'e de sons propre pays parce qu'il vantait sans cesse les avantages de la soci\'et\'e de consommation occidentale, il vit maintenant de l'aide sociale en RFA et il ne peut rien acheter hormis les choses d'usage courant.

Der Ungl\"ucklicher, nachdem er aus seinem selben Land ausgewiesen wurde, denn er immer die Vorz\"uge der Westlichen Konsumgesellschaft pries, lebt er jetzt aus Sozialhilfe des BRD, und kann sich nichts leider ausser von t\"aglichen Bedarf. 

\Question % 12. Je ne pourrai jamais me pardonner de ne m'�tre pa tu

Ich werde mir niemals verzweifeln, mich nicht schweigen zu haben.

\Question % 13. Apr\`es avoir vendu son magasin de fournitures de bureau, sa femme a ouvert une boutique de mode

Nachdem er sein B\"urobedarf Gesch\"aft verkauft hat, hat seine Frau ein Modegesch\"aft ge\"offnet.

\Question % 14. Il est soup�onn\'e d'appartenir \`a une organisation terroriste parce qu'il a occup\'e une maison vide avec quelques autres \'etudiants en sociologie.

Er steht im Verdacht, zu einer terroristischen Organisation zu geh\"oren, denn er  mit einigen Sociologie Studenten in einem leeren Haus gewohnt hat.

\end{Answer}

\end{comment}

%\begin{Exercise}[
%	title={\ldaDISK}, 
%label={sein:disk}]

%\Question % Ist der Verbraucher ein Spielzeug der Industrie?
%\Question % Der Komsumverzicht: Utopie oder Zivilcourage?
%\Question % Soll Werbung nur informieren oder auch W\"unsche wecken?
%\Question % Das heutige Konsumverhalten: ein ausgleich f\"ur die verlorene Regiosit\"at von fr\"uher?
%\end{Exercise}

\begin{comment}
\end{comment}
\chapter{\nameref{O-ch:spiel}}

\begin{Exercise}[title={\nameref{O-spiel:infsubs}}, label={spiel:infsub}]
Formez des phrases contenant un infinitif substantiv\'e
\Question der Rasen;  verboten sein; bei Strafe; betreten. 
\Question gewisse Fluglinien; sollen; rauchen; auf; verboten werden.
\Question mich; ersch�pft haben; v�llig; endlos; warten.
\Question in; sehr; baden; die meisten Fl�sse; gef�hrlich sein.
\Question eine Gesellschaft; in; ohne; zusammenleben; Spielregeln; unm�glich sein.
\Question viel trinken; seine; schaden; Gesundheit.
\Question alle Museen; in; verboten sein; mit Blitzlicht; fotografieren.
\Question Hunde; m�ssen; von; vorher; mitnehmen  angek�ndigt werden.
\Question parken; problematisch geworden; die Innenstadt; in; sein.
\Question abstellen; Fahrr�der; von; verboten sein; das Schaufenster; vor
\end{Exercise}

\begin{Answer}[ref={spiel:infsub}]
\Question Bretreten des Rasens ist bei Strafe verboten.
\Question Rauchen auf gewisse Fluglinien sollte verboten werden.
\Question Endlos Warten hat mich v\"ollig ersch\"opt.
\Question In der meisten Fl\"usse ist Baden sehr gef\"arhlich.
\Question In einer Gesellschaft ohne Spielregeln ist Zusammenleben unm\"oglich.
\Question Vieles Trinken shadet seiner Gesundheit.
\Question Mit Blitzlicht Fotografieren ist in alle Museen verboten.
\Question Das Mitnehmen von Hunde m\"uss vorher angek\"undigt werden.
\Question In der Innenstadt ist Parken problematisch geworden.
\Question Abstellen von Fahrr\"ader ist vor dem Schaufenster verboten.
\end{Answer}

\begin{Exercise}[label={spiel:infsub2}]
Traduisez les phrases que vous avez form\'ees dans l'exercice pr\'ec\'edent
\end{Exercise}

\begin{Answer}[ref={spiel:infsub2}]
\Question%1
  Marcher sur la pelouse est interdit sous peine d'amende.
\Question%2
  Fumer devrait \^etre interdit sur certaines lignes a\'eriennes.
\Question%3
  Attendre sans fin m'a \'epuis\'e.
\Question%4
  Dans la plupart des fleuves, se baigner est interdit.
\Question%5
  Dans une soci\'et\'e sans r\`egles, la cohabitation est impossible.
\Question%6
  L'exc\`es d'alcool nuit \`a la sant\'e.
\Question%7
  Prendre des photos au flash est interdite dans tous les mus\'ees.
\Question%8
  Se faire accompagner d'un chien doit \^etre annonc\'e \`a l'avance.
\Question%9
  Au centre ville, se garer est probl\'ematique.
\Question%10
  Descendre de son v\'elo est interdit devant la vitrine.
\end{Answer}

\begin{Exercise}[title={\nameref{O-spiel:wann}}, label={spiel:wann}]
\Question%1 
	Ich frage mich, \ldaBLANK ich Zeit haben werde.
\Question%2 
	\ldaBLANK er seine Forderung durchsetz, werde ich ihm gratulieren.
\Question%3 
	\ldaBLANK hat er dich betrogen?
\Question%4 
	\ldaBLANK du willst, k\"onnen wir morgen hinfahren.
\Question%5 
	Ich m\"ochte wissen, \ldaBLANK er das behauptet hat.
\Question%6 
	Ich m\"ochte gern eine Platte h\"oren, \ldaBLANK es dich nicht st\"ort.
\Question%7 
	Wir haben nicht erfahren k\"onnen, \ldaBLANK die neue Regel in Kraft tritt.
\Question%8 
	Man kann sich fragen, \ldaBLANK er aufh\"oren wird zu l\"ugen.
\Question %9
Ich bin jetz fertig: Wir k\"onnen hinfahren, \ldaBLANK du willst, sofort oder auch sp\"ater.
\Question %10
	\ldaBLANK dir die Farbe nicht gef\"allt, kannst du den Pullover umtauschen.
\end{Exercise}

\begin{Answer}[ref={spiel:wann}]
\Question%1 Ich frage mich, \ldaBLANK ich Zeit haben werde.
  Wann.
\Question%2 \ldaBLANK er seine Forderung durchsetz, werde ich ihm gratulieren.
  Wenn.
\Question%3 \ldaBLANK hat er dich betrogen?
  Wann.
\Question%4 \ldaBLANK du willst, k\"onnen wir morgen hinfahren.
  Wenn.
\Question%5 Ich m\"ochte wissen, \ldaBLANK er das behauptet hat.
  Wann.
\Question%6 Ich m\"ochte gern eine Platte h\"oren, \ldaBLANK es dich nicht st\"ort.
  Wenn.
\Question%7 Wir haben nicht erfahren k\"onnen, \ldaBLANK die neue Regel in Kraft tritt.
  Wann.
\Question%8 Man kann sich fragen, \ldaBLANK er aufh\"oren wird zu l\"ugen.
  Wann. %`ob' 
\Question%9 Ich bin jetz fertig: Wir k\"onnen hinfahren, \ldaBLANK du willst, sofort oder auch sp\"ater.
  Wann.
\Question%10 \ldaBLANK dir die Farbe nicht gef\"allt, kannst du den Pullover umtauschen.
  Wenn.
\end{Answer}

\begin{Exercise}[label={spiel:wann2}]
Traduisez les phrases de l'Exercise pr\'ec\'edent
\end{Exercise}

\begin{Answer}[ref={spiel:wann2}]
\Question%1
  Je me demande quand j'aurai le temps.
\Question%2
  Si sa d\'emarche aboutit, je le f\'eliciterai
\Question%3
  Quand t'as t-il tromp\'e.
\Question%4
  Si tu le veux, nous pouvons y aller demain.
\Question%5
  Je voudrais savoir, quand il d\'eclar\'e cela.
\Question%6
  J'aimerais bien \'ecouter un disque, si \c ca ne te d\'erange pas.
\Question%7
  Nous n'avons pas trouver, quand la nouvelle r\`egle entrera en vigueur.
\Question%8
  On peut se demander, quand il s'arr\^etera de mentir.
\Question%9
  Je suis pr\^et maintenant: nous pouvons y aller, quand tu veux, maintenant ou bien plus tard.
\Question%10
  Si la couleur ne te pla\^it pas, tu peux \'echanger le pullover contre un autre.
\end{Answer}

\begin{Exercise}[
	title={\nameref{O-spiel:simult}}, 
label={spiel:simult}]
Compl\'etez par un adverbe ou une conjonction exprimant la simutan\'eit\'e
\Question%
	Man kann doch nicht \ldaBLANK an zwei Orten sein.
\Question%2 
	\ldaBLANK du Fieber hast, darfst du nicht aus dem Haus.
\Question%3 
	\ldaBLANK du darauf bestehst, machsts du dich eher verd\"achtig.
\Question%4 
	Wir haben \ldaBLANK erfahren, dass er ein Betr\"uger ist. \ldaVERIF
\Question%5 
	Er hat es geschafft, \ldaBLANK er den zust\"andigen Beamten bestochen hat.
\Question%6 
	\ldaBLANK er seiner Freude freien Lauf l\"asst, k\"onnen seine Eltern ihre Entt\"auschung nicht vergeben.
\Question%7 
	Er wollte seinem Vater helfen und hat ihn \ldaBLANK nur gest\"ort.
\Question%8 
	\ldaBLANK er mit Freunden in der Eckkneipe sitz, sieht sie fern oder sie liest.
\Question%9 
	Wenn du die Koffer packst, kann ich \ldaBLANK ein Hotelzimmer in Heidelberg reservieren.
\Question%10 
	Er wollte die Strecke an einem Tag zur\"ucklegen und hat \ldaBLANK sein Auto kaputtgefahren.
\end{Exercise}

\begin{Answer}[ref={spiel:simult}]
\Question%1 Man kann doch nicht \ldaBLANK an zwei Orten sein.
  Gleichzeitig.
\Question%2 \ldaBLANK Fieber hast, darfst du nicht aus dem Haus.
  Solange.
\Question%3 \ldaBLANK du darauf bestehst, machsts du dich eher verd\"achtig.
  Indem.
\Question%4 Wir haben \ldaBLANK erfahren, dass er ein Betr\"uger ist.
  Gleichzeitig.
\Question%5 Er hat es geschafft, \ldaBLANK er den zust\"andigen Beamten bestochen hat.
  Indem.
\Question%6 \ldaBLANK er seiner Freude freien Lauf l\"asst, k\"onnen seine Eltern ihre Entt\"auschung nicht vergeben.
  Solange.
\Question%7 Er wollte seinem Vater helfen und hat ihn \ldaBLANK nur gest\"ort.
  Dabei. 
\Question%8 \ldaBLANK er mit Freunden in der Eckkneipe sitz, sieht sie fern oder sie liest.
  W\"ahrend.
\Question%9 Wenn du die Koffer packst, kann ich \ldaBLANK ein Hotelzimmer in Heidelberg reservieren.
  Indessen.
\Question%10 Er wollte die Strecke an einem Tag zur\"ucklegen und hat \ldaBLANK sein Auto kaputtgefahren.
  Unterdessen.
\end{Answer}

\begin{Exercise}[label={spiel:simult2}]
Traduisez les phrases de l'Exercice pr\'ec\'edent
\end{Exercise}

\begin{Answer}[ref={spiel:simult2}]
\Question%1
  On ne peut \^etre \`a deux endroits \`a la fois.
\Question%2
  Tant que du as de la fi\`evre, tu ne devras pas sortir.
\Question%3
En te comportant de la sorte, tu te rends plut\^ot suspicieux.  
\Question%4 
  Nous avons simultan\'ement r\'ealis\'e, que c'\'etait un tricheur.
\Question%5
  Il y est parvenu, en soudoyant les fonctionnaires comp\'etents.
\Question%6
  Tant qu'il laissait libre cours \`a son plaisir, ses parent ne pouvaient pas lui pardonner.
\Question%7
  Il voulait aider son p\`ere, mais il n'a r\'eussi qu\`a le d\'eranger.
\Question%8
Alors qu'il est assis au bar avec des amis, elle regarde la t\'e\'e ou lit.  
\Question%9
  Si tu fais les valises, je pourrai r\'eserver une chambre d'h\^otel \`a Heidelberg.
\Question%10
  Il voulait couvrir la distance en un jour, mais entre temps il a an\'eanti sa voiture (dans un accident).
\end{Answer}

  \begin{Exercise}[
	title={\nameref{O-spiel:wie}}, 
 label={spiel:wie}]
Traduisez en employant l'interrogatif `wie'
  \Question%1 
  	Tu ne peux pas savoir combien il est riche.
  \Question%2 
  	Combien de temps es-tu rest\'e \`a Berlin.
  \Question%3 
  	Je n'avais pas id\'ee de la gravit\'e de sa situation.  
  \Question%4 
  	Tu ne peux pas t'imaginer avec quelle facilit\'e il y est arriv\'e.
  \Question%5 
  	Quelle distance y a-t-il encore jusqu'au prochain village?
  \Question%6 
  	J'admire avec quelle prudence il a r\'epondu \`a toutes les questions.
  \Question%7 
  	Quel est le prix d'une chambre d'h\^otel avec salle de bains?
  \Question%8 
  	Quelle tristesse de le voir dans un tel \'etat!
  \Question%9 
  	Si tu savais combien il m'est d\'esagr\'eable d'avoir oubli\'e de l'inviter!    
  \Question%10 
  	Tu ne te moquerais pas de lui si tu savais comme il est scrupuleux dans tout ce qu'il fait.
  \Question%11 
  	Quel courage de sa part d'affirmer que la monarchie est un syst\`eme d\'epass\'e!
  \Question%12 
  	Quelle est la hauteur de la Tour Eiffel
  \Question%13 
  	Je peux \`a peine exprimer combien je suis heureux!
  \Question%14 
  	C'est incroyable comme il peut \^etre dangereux de dire ce que l'on pense!
  \Question%15 
  	Il voulait me montrer combien il est ridicule d'avoir mauvaise conscience pour cela.
\end{Exercise}

\begin{Answer}[ref={spiel:wie}]
  \Question%1 Tu ne peux pas savoir combien il est riche.
    Du kannst nicht wissen, wie reich er ist.
  \Question%2 Combien de temps es-tu rest\'e \`a Berlin.
    Wie lange, bist du in Berlin geblieben?
  \Question%3 Je n'avais pas id\'ee de la gravit\'e de sa situation.  
    Ich h\"atte keine Ahnung, wie schwer die Lage war.
  \Question%4 Tu ne peux pas t'imaginer avec quelle facilit\'e il y est arriv\'e.
    Tu kannst dir nicht vorstelle, wie leicht es ihm gelungen ist.
  \Question%5 Quelle distance y a-t-il encore jusqu'au prochain village?
    Wie weit vor dem n\"achsten Dorf?
  \Question%6 J'admire avec quelle prudence il a r\'epondu \`a toutes les questions.
    Ich bewundere wie vorsichtig er hat alle Fragen geantwortet.
  \Question%7 Quel est le prix d'une chambre d'h\^otel avec salle de bains?
    Wieviel kostest ein Badezimmer?
  \Question%8 Quelle tristesse de le voir dans un tel \'etat!
    Wie tra\"urig, ihn in einem solchen Zustand zu sehen.
  \Question%9 Si tu savais combien il m'est d\'esagr\'eable d'avoir oubli\'e de l'inviter!    
    Wenn du w\"usste, wie mir unangenehem ist, vergessen zu haben, ihm zu einladen?
  \Question%10 Tu ne te moquerais pas de lui si tu savais comme il est scrupuleux dans tout ce qu'il fait.
    Du w\"urdest ihn nicht auslachen, wenn du wusste, wie gewissenhaft er ist, in alles was er tut.
  \Question%11 Quel courage de sa part d'affirmer que la monarchie est un syst\`eme d\'epass\'e!
    Wie mut von ihm zu behaupten, dass die Monarchie veraltet ist.
  \Question%12 Quelle est la hauteur de la Tour Eiffel
    Wie hoch ist der Eiffelturm?
  \Question%13 Je peux \`a peine exprimer combien je suis heureux!
    Ich kann kaum ausdr\"ucken, wie glucklich ich bin.
  \Question%14 C'est incroyable comme il peut \^etre dangereux de dire ce que l'on pense!
    Es ist unglaublich, wie gef\"ahrlich es sein k\"onnte, was man denkt laut zu sagen.
  \Question%15 Il voulait me montrer combien il est ridicule d'avoir mauvaise conscience pour cela.
    Er wollte mich schauen, wie l\"acherlich es ist, deswegen ein schlechtes Gewissen zu haben.
\end{Answer}     
    
\begin{Exercise}[
	title={\ldaTHEME}, 
label={spiel:theme}]

\Question % 1 
J'ai bonne conscience d'avoir accept\'e sa proposition, car il passe en g\'en\'eral pour un honn�te homme

%Indication

\Question % 2 
	Si ma m\'emoire ne m'abuse pas, il a d\'ej� essay\'e sans succ\`es de s'immiscer dans les luttes sociales

\Question % 3 
	Depuis qu'elle s'est bless\'ee � la jambe en skiant, elle a beaucoup de mal � marcher

\Question % 4 
	J'attends avec impatience de savoir si l'agitation sociale entra�nera l'effondrement du r\'egime

\Question % 5 
	Les rapports de force au sein du parti ont chang\'e sous la pression de l'opinion publique

\Question % 6 
	Je suis plut�t d\'e�u qu'il ait chang\'e d'avis entre-temps, car les innovations qu'il voulait introduire dans le domaine social me paraissaient tr\`es importantes

%Indications: 

\Question % 7 
	Le gouvernement a d\'eclench\'e une profonde crise sociale en r\'eprimant impitoyablement l'insurrection populaire

\Question % 8 
	Alors que tous ses coll\`egues de travail partent en voyage pendant les vacances, son fr\`ere pr\'ef\`ere rester chez lui pour r\'enover sa maison

\Question % 9 
	Beaucoup de politiciens veulent se mettre en valeur en mettant l'ordre \'etabli en question

\Question % 10 
	On ne peut qu'admirer avec quelle facilit\'e il a r\'epondu � toutes les questions et combien il a \'et\'e prudent en promettant de maintenir la ligne politique de son pr\'ed\'ecesseur

\Question % 11 
	Il sait que cela ne m\`ene � rien de refuser tout compromis, mais il ne veut pas se soumettre aux exigences de l'opposition

\Question % 12 
	Pour se rendre ma�tre de la situation, le chancelier f\'ed\'eral a d� satisfaire la revendication des syndicats et \'ecarter le ministre de l'int\'erieur

\end{Exercise}    

\begin{Answer}[ref={spiel:theme}]

\Question % 1 J'ai bonne conscience d'avoir accept\'e sa proposition, car il passe en g\'en\'eral pour un honn�te homme

Ich habe gutes Gewissen, sein Vorschlag angenommen zu haben, denn er gilt als ein ehrlicher Mensch.

\Question % 2 Si ma m\'emoire ne m'abuse pas, il a d\'ej� essay\'e sans succ\`es de s'immiscer dans les luttes sociales

Wenn mich mein Ged\"atchnis  nicht tr\"ugt, % schw\"oren,\ldots
 hat er vergeblich versucht sich in sozialen einzumischen?

\Question % 3 Depuis qu'elle s'est bless\'ee � la jambe en skiant, elle a beaucoup de mal � marcher

Seitdem er sich beim Skilaufen verletzt hat, %infinitif substantiv\'e
f\"allt ihr das Laufen schwer. 

\Question % 4 J'attends avec impatience de savoir si l'agitation sociale entra�nera l'effondrement du r\'egime

Ich bin gespannt zu wissen, ob die das des Regime bringen wird.%1?2

\Question % 5 Les rapports de force au sein du parti ont chang\'e sous la pression de l'opinion publique

Ein Wechsel der Machtverh\"altnisse innerhalb der Partei ist unter dem Druck der �ffentlichen Meinung  eingetreten. % Die Gesellschaft im Wandel

\Question % 6 Je suis plut�t d\'e�u qu'il ait chang\'e d'avis entre-temps, car les innovations qu'il voulait introduire dans le domaine social me paraissaient tr\`es importantes

%Indications: 

Ich bin eher entt\"auscht, dass er inzwischen seine Meinung ge\"andert hat, % Expression de la simulatan\'eit\'e
denn die Neuerung die er im sozialen Bereich einf\"uhren wollte, galten mir als sehr wichtig. % Die Gesellschaft im Wandel

\Question % 7 Le gouvernement a d\'eclench\'e une profonde crise sociale en r\'eprimant impitoyablement l'insurrection populaire

%Die Regierung hat eine, indem er 

\Question % 8 Alors que tous ses coll\`egues de travail partent en voyage pendant les vacances, son fr\`ere pr\'ef\`ere rester chez lui pour r\'enover sa maison

\Question % 9 Beaucoup de politiciens veulent se mettre en valeur en mettant l'ordre \'etabli en question

\Question % 10 On ne peut qu'admirer avec quelle facilit\'e il a r\'epondu � toutes les questions et combien il a \'et\'e prudent en promettant de maintenir la ligne politique de son pr\'ed\'ecesseur

\Question % 11 Il sait que cela ne m\`ene � rien de refuser tout compromis, mais il ne veut pas se soumettre aux exigences de l'opposition

\Question % 12 Pour se rendre ma�tre de la situation, le chancelier f\'ed\'eral a d� satisfaire la revendication des syndicats et \'ecarter le ministre de l'int\'erieur

\end{Answer}    

%\begin{Exercise}[title={\ldaDISK}, label={spiel:disk}]
%\Question %1 Unter welchen Einfl�ssen h�rt eine Regel auf, lebendig zu sein?
%\Question %2 Ist das Wirken der sogennanten Spielverderber die unverzichtbare Voraussetzung f�r den sozialen Wandel?
%\Question %3 Haben K�nstler und Wissenschaftler - es sind die von K�stner zitierten Spielverderber - eine tiefere und langfristigere Wirkung auf die Entwicklung der Gesellschaft als die Politiker?
%\Question %4 Ist Anpasusngf�higkeit die goldene Regel des politischen Erfolgs?
%\Question %5 Was ist ein fortschrittliche Einstellung?
%\Question %6 Ist auch eine demokratische Gesellschaft auf "Spielverderber" angeweisen, um sich grundlegend zu wandeln?
%\end{Exercise}

%\begin{Answer}[ref={spiel:disk}]
%\Question %1 Unter welchen Einfl�ssen h�rt eine Regel auf, lebendig zu sein?
%\Question %2 Ist das Wirken der sogennanten Spielverderber die unverzichtbare Voraussetzung f�r den sozialen Wandel?
%\Question %3 Haben K�nstler und Wissenschaftler - es sind die von K�stner zitierten Spielverderber - eine tiefere und langfristigere Wirkung auf die Entwicklung der Gesellschaft als die Politiker?
%\Question %4 Ist Anpasusngf�higkeit die goldene Regel des politischen Erfolgs?
%\Question %5 Was ist ein fortschrittliche Einstellung?
%\Question %6 Ist auch eine demokratische Gesellschaft auf "Spielverderber" angeweisen, um sich grundlegend zu wandeln?
%\end{Answer}

\begin{comment}
\end{comment}

% ---
% Bonus
% ---
% auf etwas Verzicht leisten
% unver[s]ichtbar


%\listofexercises
%\ListOfExerciseInToc
\tableofcontents
\bibliographystyle{plainnat}
\bibliography{lang}
\end{document}
